\ProvidesPackage{csc}

\usepackage{fullpage, bookmark, fancyhdr, changepage}
\usepackage[margin=1in]{geometry}
\usepackage{amsmath, amsthm, amsfonts, amssymb, amscd, mathrsfs}
\usepackage[dvipsnames]{xcolor}
\usepackage{enumerate, enumitem}
\usepackage{graphicx, tikz-cd, forest}
\usepackage{hyperref, setspace, listings, adjustbox}
\usepackage[T1]{fontenc}
\usepackage[nomath]{lmodern}
\usepackage{textgreek}
\usepackage[normalem]{ulem}
\usepackage{wrapfig, subcaption}
\usepackage[super]{nth}
\usepackage{labelschanged}
\usepackage{algorithm}
\usepackage[noend]{algpseudocode}
\usepackage{mathpartir}

\usepackage{tabularx}

\newcommand{\ds}{\displaystyle}
\newcommand{\ttab}{\phantom{\texttt{tt}}}
\newcommand{\tab}{\hspace{1cm}}
\newcommand{\htab}{\hspace{0.5cm}}
\newcommand{\qtab}{\hspace{0.25cm}}
\newcommand{\nl}[1]{\\[#1\baselineskip]}

\newcommand{\clr}[2]{\color{#1}#2\color{black}}
\newcommand{\ul}{\uline}
\newcommand{\ol}{\overline}
\newcommand{\lf}{\left}
\newcommand{\rt}{\right}

\newcommand{\st}{\ \big|\ }
\newcommand{\Z}{\mathbb{Z}}
\newcommand{\C}{\mathbb{C}}
\newcommand{\Q}{\mathbb{Q}}
\newcommand{\R}{\mathbb{R}}
\newcommand{\N}{\mathbb{N}}
\newcommand{\E}{\mathbb{E}}
\renewcommand{\P}{\mathbb{P}}
\renewcommand{\cal}{\mathcal}
\newcommand{\ang}[1]{\langle#1\rangle}
\newcommand{\nil}{\textsc{Nil}}

\DeclareMathOperator{\lcm}{lcm}
\DeclareMathOperator{\id}{id}
\DeclareMathOperator{\Id}{Id}
\let\Re\relax
\DeclareMathOperator{\Re}{Re}
\let\Im\relax
\DeclareMathOperator{\Im}{Im}
\let\oldhash\#%
\DeclareRobustCommand{\#}{\adjustbox{valign=B,totalheight=.57\baselineskip}{\oldhash}}

\renewcommand{\thesection}{Part \arabic{section}: \hspace{-1em}}
\renewcommand{\thesubsection}{(\alph{subsection})}

\newtheorem{theorem}{Theorem}
\newtheorem*{theorem*}{Theorem}
\newtheorem{cor}{Corollary}
\newtheorem*{cor*}{Corollary}
\newtheorem{claim}{Claim}
\newtheorem*{claim*}{Claim}
\newtheorem{lemma}{Lemma}
\newtheorem*{lemma*}{Lemma}

\theoremstyle{definition}
\newtheorem{defn}{Defn}
\newtheorem*{defn*}{Defn}

\hypersetup{
  colorlinks=true,
  linkcolor=blue,
  linkbordercolor={0 0 1}
}
 
\setlength{\footskip}{6em}
\usepackage{indentfirst}

\newcommand{\todo}[1]{\textcolor{red}{\textbf{TODO:} #1}}

\definecolor{darkblue}{rgb}{0,0,0.5}
\definecolor{darkgreen}{rgb}{0,0.3,0}
\definecolor{darkpink}{rgb}{0.4,0,0.3}
\definecolor{graygreen}{rgb}{0.3,0.5,0.3}
\definecolor{grayblue}{rgb}{0.2,0.2,0.6}
\definecolor{grayred}{rgb}{0.5,0.2,0.2}

\lstset{
  backgroundcolor=\color{white},     % choose the background color; you must add \usepackage{color} or \usepackage{xcolor}; should come as last argument
  % identifierstyle=\color{red},
  basicstyle=\normalsize\ttfamily,      % the size of the fonts that are used for the code
  breakatwhitespace=false,              % sets if automatic breaks should only happen at whitespace
  breaklines=false,                     % sets automatic line breaking
  captionpos=b,                         % sets the caption-position to bottom
  abovecaptionskip=-3 mm,
  commentstyle=\itshape\color{graygreen}, % comment style
  % escapeinside={(:}{:)},             % if you want to add LaTeX within your code
  escapechar={!},
  % extendedchars=true,              % lets you use non-ASCII characters; for 8-bits encodings only, does not work with UTF-8
  % firstnumber=1000,                % start line enumeration with line 1000
  % frame=tb,                        % adds a frame around the code
  % keepspaces=true,                 % keeps spaces in text, useful for keeping indentation of code (possibly needs columns=flexible)
  keywordstyle=\color{darkblue},     % keyword style
  language=Haskell,                  % the language of the code
  morekeywords={ Set, Tree, Leaf, Node, Applicative, fmap, liftA2, bimap, foldMap
               , traverse, mappend, pure, Foldable, Traversable, zero, one
               , Semiring, Semigroup, NonEmpty, sconcat, TSet,
               , SimplicialSet, TSimplicialSet, Graph, TGraph, LGraph
               , Map, IsString, fromString },
  deletekeywords={instance, data, where, class, filter, type, insert, delete, union, map},      % if you want to delete keywords from the given language
  emph={data, class, instance, where, type},
  emphstyle=\color{darkpink},
  numbers=none,                      % where to put the line-numbers; possible values are (none, left, right)
  literate=*{lam}{{$\lambda$}}1  {->}{{$\rightarrow$}}1 % {\\}{{$\lambda$}}1 
  {|->}{{$\mapsto$}}1 {o+o}{{$\oplus$}}1 {=>}{{$\Rightarrow$}}1
  {/\\}{{$\Lambda$}}1  {'\{}{'\{}1 {\\/}{{$\forall$}}1 %
  ,
  % numbersep=5pt,                   % how far the line-numbers are from the code
  % numberstyle=\tiny\color{mygray}, % the style that is used for the line-numbers
  % rulecolor=\color{black},         % if not set, the frame-color may be changed on line-breaks within not-black text (e.g. comments (green here))
  % showspaces=false,                % show spaces everywhere adding particular underscores; it overrides 'showstringspaces'
  % showstringspaces=false,          % underline spaces within strings only
  % showtabs=false,                  % show tabs within strings adding particular underscores
  % stepnumber=2,                    % the step between two line-numbers. If it's 1, each line will be numbered
  stringstyle=\color{grayred},     % string literal style
  % tabsize=2,                       % sets default tabsize to 2 spaces
  % title=\lstname                   % show the filename of files included with \lstinputlisting; also try caption instead of title
  xleftmargin=10pt,
  aboveskip=8pt,
  belowskip=4pt
}

\newcommand{\code}[1]{\lstinline[mathescape]|#1|}
\newcommand{\hcode}[1]{{\color{darkblue} \lstinline[keywordstyle={}]|#1|}} % h for "highlighted"
\newcommand{\hfcode}[1]{{\color{darkblue} \lstinline[keywordstyle={},basicstyle=\footnotesize\ttfamily]|#1|}} % h for "highlighted"
